%KI-Verzeichnis / AI Directory
%Required by the HTW academic integrity declaration ("KI-kennzeichnungspflichtig").
%Structure follows the HTW Handreichung (Tool / Form of use / Reason / Affected
%parts / Notes), extended with representative prompts as required by the
%declaration text.
%
%TODO before submission: check every chapter/section reference in the
%"Affected parts" column and delete any row that does not apply.

\section{AI Directory}
\label{sec:ai-directory}

The following directory lists all AI-based tools used in the preparation of this
thesis, together with the form and purpose of their use and the parts of the work
affected. Representative prompts are given in Section~\ref{subsec:ai-prompts}. All
content generated with these tools was checked against the primary literature and
against our own experimental data before it was adopted. Numerical results,
measurements and their interpretation originate from our own laboratory work.

\begin{table}[H]
\centering
\caption{Directory of AI-based tools used in this thesis}
\label{tab:ai-directory}
\footnotesize
\renewcommand{\arraystretch}{1.25}
\begin{tabularx}{\textwidth}{L{2.2cm} L{3.0cm} X L{2.6cm}}
\toprule
\thd{Tool} & \thd{Form of use} & \thd{Reason} & \thd{Affected parts} \\
\midrule

DeepL \newline (DeepL SE) &
Translation of source material &
Reading German-language literature, standards and equipment documentation in
English &
Chapters 2 and 3 \\
\addlinespace

&
Translation of own drafts &
Transferring passages drafted in German or Indonesian into the English
manuscript &
Chapters 1--5 \\
\addlinespace

&
Rephrasing (DeepL Write) &
Grammar and wording correction of sentences written by the authors &
Chapters 1--5 \\
\addlinespace
\midrule

ChatGPT \newline (OpenAI) &
Explanation of concepts &
Building an understanding of pyrolysis, steam activation and the analytical
methods as background for our own formulations &
Chapters 2 and 3 \\
\addlinespace

&
Language editing &
Grammar, style and consistency of terminology in text written by the authors &
Chapters 1--5 \\
\addlinespace

&
Draft formulations &
Producing first drafts of individual passages, which were subsequently revised,
verified and rewritten by the authors &
Chapters 2, 4 and 5 \\
\addlinespace

&
\LaTeX{} and data handling &
Assistance with \LaTeX{} code, figure and table formatting, and spreadsheet
calculations on our own measurement data &
Chapters 4 and 5, Appendix \\
\addlinespace
\midrule

Claude \newline (Anthropic) &
Drafting and revision of chapter text &
Structuring and formulating sections on the basis of our own results and notes;
all output was reviewed and edited by the authors &
Chapters 1--5 \\
\addlinespace

&
\LaTeX{} implementation &
Preparation of \texttt{.tex} source, pgfplots figures, tables and the uniform
figure and table style &
Chapters 4 and 5, Appendix \\
\addlinespace

&
Discussion of own data &
Working through the experimental results and discussing possible interpretations;
the measurements and the final interpretation are our own &
Chapter 4 \\
\addlinespace

&
Literature notes and consistency check &
Summarising literature we had selected, and checking the manuscript for internal
consistency, cross-references and language errors &
Chapters 1--5 \\
\addlinespace
\midrule

Google Gemini &
Literature search &
Identifying and pre-sorting publications, which were then retrieved and read in
the original &
Chapters 2 and 4 \\
\addlinespace

&
Explanation of concepts &
Clarifying individual technical relationships as background reading &
Chapters 2 and 3 \\

\bottomrule
\end{tabularx}
\end{table}

\vspace{1em}

\subsection{Representative prompts}
\label{subsec:ai-prompts}

The prompts below are representative examples of the inputs used for each form of
use listed in Table~\ref{tab:ai-directory}. They are reproduced in the form in
which they were entered.

\vspace{0.5em}
\noindent\textbf{DeepL}
\begin{itemize}
	\setlength{\itemsep}{0.2em}
	\item Translation of German source passages into English (text input, no
	instruction prompt).
	\item DeepL Write: rephrasing of own English sentences for grammar and style.
\end{itemize}

\noindent\textbf{ChatGPT}
\begin{itemize}
	\setlength{\itemsep}{0.2em}
	\item ``Explain the difference between physical activation with steam and with
	CO$_2$, and why steam activation is usually carried out at lower
	temperatures.''
	\item ``Correct the grammar and improve the readability of the following
	paragraph without changing its technical content: [own text].''
	\item ``How is the Van Krevelen diagram interpreted for biochars, and what do
	low H/C and O/C ratios indicate?''
	\item ``Write a LaTeX table with the following columns and values: [own
	measurement data].''
\end{itemize}

\noindent\textbf{Claude}
\begin{itemize}
	\setlength{\itemsep}{0.2em}
	\item ``Draft the results section for the H$_2$O activation experiments using
	the following mass-loss data and my notes: [own data and notes].''
	\item ``Create a pgfplots figure of biochar yield against pyrolysis temperature
	from these values: [own data].''
	\item ``Given these CHNS results, which trends are worth discussing, and how do
	they compare with the values reported in the cited literature?''
	\item ``Check the whole thesis for inconsistencies between chapters, incorrect
	cross-references and language errors, and list the findings.''
\end{itemize}

\noindent\textbf{Google Gemini}
\begin{itemize}
	\setlength{\itemsep}{0.2em}
	\item ``Which publications deal with steam activation of digestate-derived
	biochar and its use for H$_2$S removal?''
	\item ``Explain the mechanism of H$_2$S partial oxidation on a carbon surface.''
\end{itemize}
